% preamble/floats.tex — figures, tables, code listings, plots.
%
% ACCESSIBILITY RULE: every picture needs a text alternative.
%   \includegraphics[width=..., alt={What the image shows}]{file}
%   \begin{tikzpicture}[alt={What the diagram shows}] ... \end{tikzpicture}
% Purely decorative graphics use [artifact] instead of alt.

\usepackage{graphicx}
\graphicspath{{figures/}}

% \addimg{file}{width}{caption}{label}{alt text}: a complete figure with an
% image whose alt text cannot be forgotten (it is a required argument).
% Example: \addimg{sample-photo}{0.5\linewidth}{Test rig}{fig:rig}
%            {Photo of the test rig: two servers in a rack}
\NewDocumentCommand\addimg{m m m m m}{%
  \begin{figure}[htbp]
    \centering
    \includegraphics[width=#2, alt={#5}]{#1}
    \caption{#3}\label{#4}
  \end{figure}%
}

% Number figures, tables and equations per section: Figure 3.2, (2.1).
\counterwithin{figure}{section}
\counterwithin{table}{section}
\counterwithin{equation}{section}

% The float package defines the extra "Listing" float type below.
% Note: its [H] placement is not tagged correctly yet; use [htbp].
\usepackage{float}

% Tables: booktabs rules and siunitx number alignment (both compatible).
\usepackage{booktabs}
\usepackage{siunitx}
\sisetup{group-minimum-digits=4}
% Table header cells: the first row of every table is tagged TH (header),
% so screen readers can announce the column for each cell. A table with two
% header rows sets \tagpdfsetup{table/header-rows={1,2}} inside its table
% environment (see tab:performance in chapters/06-experiments.tex).
\tagpdfsetup{table/header-rows={1}}

% Diagrams and data plots. Plots read their numbers from figures/data/*.csv,
% so the figure is regenerated whenever the data changes.
\usepackage{tikz}
\usetikzlibrary{arrows.meta, positioning, shapes.geometric, calc}
\usepackage{pgfplots}
\pgfplotsset{compat=1.18, compat/show suggested version=false}

% Code listings. The listings package is not yet compatible with tagging in
% TeX Live 2026 (it breaks the structure tree), so code is typeset with the
% kernel's verbatim environment inside a numbered "Listing" float.
% Keep code lines under ~70 characters.
% Captions go above listings, as for tables (plaintop).
\floatstyle{plaintop}
\newfloat{codelisting}{htbp}{lol}[section]
\floatname{codelisting}{Listing}
\AddToHook{env/verbatim/begin}{\footnotesize\setstretch{1}}
